% !Mode:: "TeX:UTF-8"
\begin{conclusions}

视觉感知模型从训练条件走向目标场景时，既可能遭遇数据分布与类别词表的变化，也可能面对更高的任务语义要求。本文将模型在这些条件变化下仍需保持相应能力的情形统称为多泛化场景，沿域泛化、类泛化与语义理解泛化三个层面展开研究。研究对象由跨分布的样本表征、开放类别的概念识别，逐步拓展至高阶语义解释：前两个层面通过目标域适应与跨模型先验融合改善迁移表现，第三个层面则在不更新模型的条件下，以系统评测考察已有能力面对艺术情感推理任务时的适用范围。能力提升与能力评测由此构成本文两种不同而相互衔接的研究方式。

% 本文以视觉模型迁移为主线，沿域泛化、类泛化与语义理解泛化三个层面协同展开研究：域泛化层面贯穿第 2 章与第 3 章，分别沿卷积神经网络与状态空间模型两条骨干路径协同抑制数据域迁移不确定性，并进一步扩展至稀缺标注设定下的少监督迁移问题；类泛化层面对应第 4 章，面向开放词汇分割任务设计跨模型先验融合方法；语义理解泛化层面对应第 5 章，构建艺术情感推理的多层级评测协议以系统化诊断现有模型的语义理解能力。

（1）面向无监督跨域迁移，源模型依赖的域特异性线索可能使目标域表征与预测偏移；错误预测一旦被用作伪标签，又会在自训练中放大初始偏差。针对监督信号的可靠性，本文提出 UPA，利用特征空间的邻域一致性估计样本级伪标签可靠性，自适应筛选训练样本，并通过类别感知对比学习抑制噪声传播，无需引入额外的不确定性估计模型。针对骨干表征的稳定性，本文提出 SfMamba，以通道级视觉状态空间模块增强跨通道建模，并通过语义一致性打乱策略削弱背景空间共现带来的虚假相关性，在保持线性序列复杂度的同时建模长程依赖。两项工作分别从伪标签监督和特征表征环节减轻迁移误差的累积，在 Office、Office-Home、VisDA-C 与 DomainNet-126 等无源域适应基准上验证了有效性。

（2）面向少监督跨域迁移，有限标注能否有效覆盖目标域差异、并转化为目标模型的判别能力，是提高监督信息利用效率的关键。本文以固定监督预算下的目标域学习收益界定信息最大化目标，在无源主动域适应设定中提出 DAMF：通过双聚焦样本选择获取具有代表性和判别价值的目标样本，适配视觉—语言模型以提供互补语义监督，再通过交替蒸馏促进其与源模型之间的双向知识传递。协作风险上界从样本覆盖、模型风险与跨模型对齐等方面分析了方法的作用。在 Office-Home、VisDA-C 与 Office 的闭集和开集实验中，该方法相较所比较的基线取得了性能提升，表明有限标注与多模态先验的协同利用能够改善目标域适应。

（3）面向类泛化中的开放词汇分割任务，模型既要识别训练阶段未见的类别，也要准确确定其像素区域。现有方法往往难以兼顾视觉—语言模型的开放语义能力与分割基础模型的空间定位能力。本文提出 CMPF，在轻量解码器中通过查询注入和特征注入融合两类先验，并在推理阶段集成分割基础模型的零样本掩码预测，缓解类别判别与掩码质量之间的错配。该方法在全景、语义和实例分割等八个基准上验证了对未见类别的泛化能力。

（4）面向语义理解要求提升的艺术情感推理任务，已有研究分别提供了情感解释、视觉属性、刺激区域及多任务情感评测，但在同一评测体系中追踪视觉线索、情感判断与证据化解释之间的关系仍有不足。本文构建 AER-Bench，将艺术情感理解划分为感知维度、情感与效价—唤醒度—支配度、基于证据的推理和情感描述四个互补层级，结合单图与配对任务、确定性指标与开放式评分，对 14 个代表性多模态大语言模型进行系统评测。结果揭示模型的情感判断与证据化解释之间仍存在能力差距。本章通过评测刻画已有能力在高阶语义任务中的适用范围，并不将评测结果视作模型优化后的迁移增益。

综合上述工作，本文从跨分布条件下的可靠适应出发，进一步研究开放类别下的先验融合，并以高阶语义任务中的能力评测检验模型已有知识的适用范围。这一路径将稳定表征、未见类别识别与证据化语义解释置于同一目标条件变化的框架下，也明确区分了迁移方法带来的性能改进与评测所揭示的能力边界。

本文的实验主要在预设的目标域、类别集合与评测任务上展开，尚未覆盖目标条件持续变化的长期部署过程。后续研究可进一步考察分布漂移与新类别出现时的持续适应，兼顾新知识获取和已有能力保持；在标注稀缺、语义判断存在主观差异的任务中，建立能够区分模型失误、标注分歧与评测偏差的评价机制；并探索隐私与合规约束下的分布式知识迁移，使模型能够在数据无法集中汇集时仍获得可靠的适应能力。

\end{conclusions}
